% ============================================================
%  EEM 441 Control & Instrumentation / FPGA Lab Manual
%  Shared preamble -- fonts, colors, boxes, code style, TikZ
%  Compile with: lualatex -shell-escape
% ============================================================

\usepackage{fontspec}
\usepackage{unicode-math}

% ---------- Fonts ----------
\setmainfont{IBMPlexSans}[
  Path            = fonts/PlexSans/,
  Extension       = .ttf,
  UprightFont     = IBMPlexSans-Regular,
  ItalicFont      = IBMPlexSans-Italic,
  BoldFont        = IBMPlexSans-SemiBold,
  BoldItalicFont  = IBMPlexSans-SemiBoldItalic,
]

\newfontfamily\plexlight{IBMPlexSans}[
  Path            = fonts/PlexSans/,
  Extension       = .ttf,
  UprightFont     = IBMPlexSans-Light,
  ItalicFont      = IBMPlexSans-LightItalic,
]

\newfontfamily\plexmedium{IBMPlexSans}[
  Path            = fonts/PlexSans/,
  Extension       = .ttf,
  UprightFont     = IBMPlexSans-Medium,
  ItalicFont      = IBMPlexSans-MediumItalic,
]

\setmonofont{JetBrainsMono}[
  Path            = fonts/JetBrainsMono/,
  Extension       = .ttf,
  UprightFont     = JetBrainsMono-Regular,
  ItalicFont      = JetBrainsMono-Italic,
  BoldFont        = JetBrainsMono-Bold,
  BoldItalicFont  = JetBrainsMono-BoldItalic,
  RawFeature      = {-calt},   % disable ligatures: <= /= -> := must stay literal
  Scale           = MatchLowercase,
]

\setmathfont{STIXTwoMath-Regular}[
  Path      = fonts/STIX2/,
  Extension = .ttf,
]

% ---------- Colour palette ----------
\usepackage{xcolor}
\definecolor{usmpurple}{HTML}{4B1D6B}   % USM crest purple, used sparingly for accents
\definecolor{inkblack}{HTML}{1A1A1A}
\definecolor{slategray}{HTML}{5B6570}
\definecolor{panelbg}{HTML}{F7F7F9}
\definecolor{panelrule}{HTML}{D8D8DE}
\definecolor{accentblue}{HTML}{2A5DB0}
\definecolor{accentteal}{HTML}{1B7A72}
\definecolor{warnbg}{HTML}{FFF6E5}
\definecolor{warnrule}{HTML}{E8B93F}
\definecolor{codebg}{HTML}{F5F6FA}

\color{inkblack}

% ---------- Page geometry ----------
\usepackage[a4paper,margin=2.2cm,headheight=28pt]{geometry}

% ---------- Headers / footers ----------
\usepackage{fancyhdr}
\pagestyle{fancy}
\fancyhf{}
\renewcommand{\headrulewidth}{0.4pt}
\renewcommand{\footrulewidth}{0pt}
\fancyhead[L]{\footnotesize\plexmedium\color{slategray}\leftmark}
\fancyhead[R]{\footnotesize\plexmedium\color{slategray}EEM 441 --- Control, Instrumentation \& FPGA Lab}
\fancyfoot[C]{\small\color{slategray}\thepage}

% ---------- Section styling ----------
\usepackage{titlesec}
\titleformat{\section}
  {\normalfont\Large\plexmedium\color{inkblack}}
  {\thesection}{0.6em}{}
\titleformat{\subsection}
  {\normalfont\large\plexmedium\color{inkblack}}
  {\thesubsection}{0.6em}{}
\titlespacing*{\section}{0pt}{1.4em}{0.8em}
\titlespacing*{\subsection}{0pt}{1.1em}{0.6em}

% ---------- Lists ----------
\usepackage{enumitem}
\setlist{topsep=4pt, itemsep=2pt, parsep=0pt}

% ---------- Tables ----------
\usepackage{array}
\usepackage{booktabs}
\usepackage{longtable}
\usepackage{multirow}
\usepackage{caption}
\captionsetup{font={small}, labelfont={color=accentblue}, skip=6pt}

% Non-floating table/figure wrappers -- required inside tcolorbox environments,
% since LaTeX floats (table/figure) cannot be used inside non-page boxes.
\newenvironment{boxedtable}[1][]{%
  \par\medskip\centering
}{%
  \par\medskip
}
\newenvironment{boxedfigure}[1][]{%
  \par\medskip\centering
}{%
  \par\medskip
}

% ---------- Hyperlinks ----------
\usepackage[colorlinks=true, linkcolor=accentblue, urlcolor=accentblue, citecolor=accentblue]{hyperref}

% ---------- Graphics / TikZ ----------
\usepackage{graphicx}
\usepackage{tikz}
\usetikzlibrary{shapes.geometric, arrows.meta, positioning, calc, fit, backgrounds}

\tikzset{
  block/.style = {
    draw=slategray, thick, fill=panelbg, rounded corners=1pt,
    minimum height=1.1cm, minimum width=2.6cm, align=center,
    font=\small\plexmedium
  },
  io/.style = {
    draw=accentblue, thick, fill=white, rounded corners=1pt,
    minimum height=0.9cm, minimum width=2.2cm, align=center,
    font=\small
  },
  decision/.style = {
    draw=slategray, thick, fill=panelbg, diamond, aspect=2,
    align=center, font=\small\plexmedium, inner sep=2pt
  },
  startstop/.style = {
    draw=accentteal, thick, fill=white, rounded corners=10pt,
    minimum height=0.8cm, minimum width=2cm, align=center,
    font=\small\plexmedium
  },
  arr/.style = {-{Latex[length=2.2mm]}, thick, draw=slategray},
}

% ---------- Minted / Pygments ----------
\usepackage{minted}
\setminted{
  fontsize     = \small,
  breaklines   = true,
  breakanywhere= false,
  linenos      = true,
  numbersep    = 6pt,
  frame        = none,
  tabsize      = 2,
  baselinestretch = 1.05,
}
\setmintedinline{fontsize=\small}
\usemintedstyle{friendly}

\renewcommand{\theFancyVerbLine}{\small\color{slategray}\arabic{FancyVerbLine}}

% Wrap minted in a light panel consistent with the section boxes
\newenvironment{codepanel}[1]{%
  \begin{tcolorbox}[
      enhanced, breakable,
      colback=codebg, colframe=panelrule,
      boxrule=0.6pt, arc=1.5pt,
      left=4pt, right=4pt, top=3pt, bottom=3pt,
      title={\small\plexmedium\color{slategray} #1},
      fonttitle=\small, coltitle=slategray,
      colbacktitle=panelbg, boxsep=1pt,
    ]
}{%
  \end{tcolorbox}
}

% ---------- Section panels (Objectives / Equipment / Procedure / etc.) ----------
\usepackage[most]{tcolorbox}

\newtcolorbox{labsection}[2]{
  enhanced, breakable,
  colback=panelbg, colframe=panelrule,
  boxrule=0.6pt, arc=2pt,
  left=8pt, right=8pt, top=6pt, bottom=6pt,
  title={\plexmedium\color{inkblack} #1},
  fonttitle=\normalsize,
  attach title to upper=\par\vspace{2pt}\noindent,
  coltitle=inkblack,
  colbacktitle=panelbg,
  toptitle=2pt, bottomtitle=2pt,
  overlay={
    \draw[accentblue, line width=2.2pt]
      ([xshift=0pt]frame.north west) -- ([xshift=0pt]frame.south west);
  },
  #2
}

% Convenience wrappers for the standard recurring sections
\newenvironment{objectives}{\begin{labsection}{Objectives}{}}{\end{labsection}}
\newenvironment{equipment}{\begin{labsection}{Equipment / Software Required}{}}{\end{labsection}}
\newenvironment{introduction}{\begin{labsection}{Introduction}{}}{\end{labsection}}
\newenvironment{procedure}{\begin{labsection}{Procedure}{}}{\end{labsection}}
\newenvironment{materials}{\begin{labsection}{Materials}{}}{\end{labsection}}
\newenvironment{reportq}{\begin{labsection}{Report}{}}{\end{labsection}}

% ---------- Placeholder / TODO callout (for hardware not yet finalized) ----------
\newtcolorbox{needsupdate}{
  enhanced, breakable,
  colback=warnbg, colframe=warnrule,
  boxrule=0.7pt, arc=2pt,
  left=8pt, right=8pt, top=5pt, bottom=5pt,
  fonttitle=\small\plexmedium,
  title={\color{warnrule!70!black} Pending},
  coltitle=warnrule!60!black,
  colbacktitle=warnbg,
}

% ---------- Experiment title block ----------
\newcommand{\experimentheader}[3]{%
  % #1 = experiment number, #2 = short title, #3 = full title
  \clearpage
  \phantomsection
  \addcontentsline{toc}{section}{Experiment #1: #2}
  \markboth{Experiment #1 --- #2}{Experiment #1 --- #2}
  \begingroup
  \centering
  \vspace*{0.4em}
  {\plexlight\color{slategray}\large EEM 441 --- Control, Instrumentation \& FPGA Laboratory\par}
  \vspace{0.6em}
  {\plexmedium\color{usmpurple}\Huge Experiment #1\par}
  \vspace{0.2em}
  {\plexmedium\color{inkblack}\LARGE #3\par}
  \vspace{0.8em}
  \endgroup
}

% ---------- Table of contents styling ----------
\usepackage{titletoc}
\renewcommand{\contentsname}{\plexmedium\color{inkblack}Table of Contents}
\titlecontents{section}
  [0pt]
  {\vspace{6pt}\plexmedium}
  {}{}
  {\titlerule*[6pt]{.}\contentspage}

% ---------- Misc ----------
\usepackage[english]{babel}
\usepackage{csquotes}
\frenchspacing
\setlength{\parindent}{0pt}
\setlength{\parskip}{6pt}
